\scanpage{125}\mspage{70}
J'ai envie de grouper ce qui a été prouvé dans un énoncé récapitulatif.

\uline{Théorème}\inserted{-\uline{Scholie}} \quad On considère le diagramme de catégories de la page 8.\textsuperscript{*}
% slanted
\margin{On suppose sur $W$ les conditions a) b) c) de p.\,\unsure{4 pr.|4 sq.|4}, plus L 5 (p.\,\unsure{43}). Cf. Commentaires p.\,\unsure{81}-\deleted{\unsure{82}}}

\textcircled{I} Quand on munit toutes les 13 catégories du diagramme du localiseur $W_S$ ($W$-équivalences ``grossières'', i.e. $W$-équivalence de $\mathrm{Cat}$, sans plus), \uline{tou\deleted{\unsure{tes|te}}s} les foncteurs du diagramme induisent des équivalences sur les catégories localisées.\textsuperscript{*}
% slanted
\margin{De plus, si $\unsure{\mathcal F|F}$ est une sous-catégorie pleine de $\mathrm{Cat}_{/S}$ qui contient une des deux sous-catégories $\text{\uline{Fib}}_0(\unsure{S})$, $\text{\uline{Cofib}}_0 S$, \deleted{$\unsure{\mathrm{Parf}\,S}$} (ce qui inclut le cas de 8 catégories du diagramme, autres que $\mathrm{Cat}_{/S}$ \deleted{\unsure{elle}}-même), le foncteur $(\mathcal F,W_S\mid\mathcal F)$ dans $(\mathrm{Cat}_{/S},W_S)$ induit une équiv. pour les localisées.}
Désignons par
\[
 \mathrm{HOT}^{\unsure{\mathrm{lc}|\ell\mathrm{c}}}_W(S)
\]
``la'' catégorie localisée commune.

\textcircled{II} Munissons les quatre catégories \deleted{\unsure{bordantes}} gauches
\[
 \mathrm{Cat}_{/S}\hookleftarrow\text{\uline{Liss}}\,S\hookleftarrow\text{\uline{Fib}}\,S\longleftarrow\text{\uline{Fibsc}}\,S
 \tag*{$(*)$}
\]
du localiseur $W_S^d$ de la $W$-équivalence fine \uline{droite} sur $S$ (
\diagram{\begin{tikzcd}[ampersand replacement=\&,column sep=tiny,row sep=tiny]
 f:X\arrow[rr]\arrow[dr,no head]\&{}\&Y\arrow[dl,no head]\\
 {}\&S\&{}
\end{tikzcd}}
est dans $W_S^d$ ssi les ${}_{s\backslash}X\unsure{\longrightarrow|\text{---}}{}_{s\backslash}Y$ le sont, pour tout $s\in\mathrm{Ob}\,S$). Les foncteurs du diagramme i.e. ceux de $(*)$ entre ces catégories, \deleted{\unsure{et}} \inserted{donc aussi} leurs composés, induisent des équivalences entre les localisées.\textsuperscript{*}
% slanted
\margin{Plus généralement, pour toute sous-catégorie pleine $\mathcal F$ de $\mathrm{Cat}_{/S}$ qui contient $\text{\uline{Fib}}\,S$ (et dualement, avec $W_S^g$ au lieu de $W_S^d$, si $\mathcal F$ contient $\text{\uline{Cofib}}\,S$).}
Désignons par
\[
 \mathrm{HOT}_W(S)
\]
``la'' catégorie localisée commune.

Enoncé dual pour
\[
 \mathrm{Cat}_{/S}\hookleftarrow\text{\uline{Prop}}\,S\hookleftarrow\text{\uline{Cofib}}\,S\hookleftarrow\text{\uline{Cofibsc}}\,S,
 \tag*{$(*')$}
\]
et les localiseurs $W_S^g$ dans ces catégories ($f\in W_S^g$ ssi les $f_{/s}:X_{/s}\longrightarrow Y_{/s}$ sont dans $W$). On trouve une ``même catégorie localisée, notée
\[
 \mathrm{HOT}_W^\circ(S),
\]
et on \deleted{\unsure{a}} \deleted{\unsure{un}} tautologiquement
% archive sheet number: 124
